\section{Alternative multipole fitting range $30 \leq \ell \leq 199$}
\label{sec:appendix_2}

This appendix presents the results of the global power-law fitting analysis carried out over the multipole range $30 \leq \ell \leq 199$, as a robustness check against the baseline range $20 \leq \ell \leq 199$ adopted in Section~\ref{sec:results}. Starting at $\ell = 30$ is the conventional choice in the literature \citep{QUIJOTE_IV, Martire_2022}, as it avoids the largest angular scales where residual systematic effects and foreground modelling uncertainties may be more significant.

However, as discussed in Section~\ref{sec:results}, excluding the first multipole bin ($20 \leq \ell \leq 39$, $\ell_{\rm eff} = 29$) has a marked impact on the BB-mode power-law fit. The BB polarization power spectrum is intrinsically fainter than EE and falls steeply with multipole; as a result, the first bin carries a substantial fraction of the total constraining power on the BB angular spectral index $\alpha_s^{BB}$. Without it, the remaining bins span a narrower dynamic range in signal amplitude, and the power-law slope becomes poorly constrained. This manifests as an artificial steepening of the inferred $\alpha_s^{BB}$ towards unphysically large negative values, accompanied by significantly enlarged posterior uncertainties. The effect is most pronounced in the $|b| > 10^\circ$ mask when QUIJOTE data are included, where the additional low-frequency sensitivity of QUIJOTE makes the model sensitive to the shape of the synchrotron angular spectrum, amplifying the impact of missing low-$\ell$ information. The inclusion of the first multipole bin in the baseline analysis resolves this degeneracy, yielding well-constrained and physically consistent estimates for both the EE and BB modes.

The results below are presented for both the auto+cross-spectra and cross-spectra-only configurations, and for the three Galactic masks, to allow a direct comparison with the baseline results of Section~\ref{sec:results} and Appendix~\ref{sec:appendix_1}.

\subsection{Fitting all data (auto- and cross-spectra)}

\textcolor{red}{RESULTS AUTO+CROSS-SPECTRA -- $30 \leq \ell \leq 199$}


\begin{table*}[h]
\centering
\footnotesize
\setlength{\tabcolsep}{4pt}
\caption{Posterior constraints (median and $68\%$ credible intervals) on the synchrotron and dust amplitudes, spectral indices, and synchrotron-dust correlation coefficient $\rho$, from a global power-law fit to the EE and BB polarization auto- and cross-spectra of the WP+Pl and QJ+WP+Pl datasets over the multipole range $30 \leq \ell \leq 200$, for the three Galactic masks.}
\label{tab:fit_all_masks_all_pairs_ell30-200}
\begin{adjustbox}{center}
\scalebox{1}{
\begin{tabular}{llllcccccccc}
\toprule
Mask & Dataset & Mode & $A_s\,[10^{-3}\,\mu\mathrm{K}^2]$ & $\alpha_s$ & $\beta_s$ & $A_d\,[10^{-3}\,\mu\mathrm{K}^2]$ & $\alpha_d$ & $\beta_d$ & $\rho$ \\
\midrule
\multirow{4}{*}{$|b|>10^\circ$}
 & WP+Pl    & EE & $9.451^{+0.649}_{-0.644}$ & $-2.845^{+0.146}_{-0.156}$ & $-3.189^{+0.116}_{-0.120}$ & $2.999^{+0.012}_{-0.012}$ & $-2.244^{+0.006}_{-0.006}$ & $1.523^{+0.006}_{-0.006}$ & $0.090^{+0.011}_{-0.011}$ \\
 &          & BB & $2.589^{+0.534}_{-0.530}$ & $-2.380^{+0.443}_{-0.506}$ & $-3.402^{+0.295}_{-0.325}$ & $2.336^{+0.008}_{-0.008}$ & $-2.163^{+0.006}_{-0.006}$ & $1.536^{+0.005}_{-0.005}$ & $0.125^{+0.020}_{-0.019}$ \\
 & QJ+WP+Pl & EE & $8.806^{+0.325}_{-0.328}$ & $-2.949^{+0.079}_{-0.083}$ & $-3.148^{+0.036}_{-0.036}$ & $2.999^{+0.012}_{-0.012}$ & $-2.244^{+0.006}_{-0.006}$ & $1.523^{+0.006}_{-0.006}$ & $0.088^{+0.008}_{-0.008}$ \\
 &          & BB & $1.266^{+0.296}_{-0.284}$ & $-4.074^{+0.449}_{-0.535}$ & $-3.065^{+0.109}_{-0.110}$ & $2.336^{+0.008}_{-0.008}$ & $-2.163^{+0.006}_{-0.006}$ & $1.535^{+0.005}_{-0.005}$ & $0.103^{+0.016}_{-0.015}$ \\
\midrule
\multirow{4}{*}{$|b|>15^\circ$}
 & WP+Pl    & EE & $8.697^{+0.678}_{-0.677}$ & $-2.695^{+0.169}_{-0.181}$ & $-3.234^{+0.138}_{-0.143}$ & $1.889^{+0.011}_{-0.011}$ & $-2.112^{+0.009}_{-0.009}$ & $1.533^{+0.008}_{-0.008}$ & $0.066^{+0.012}_{-0.012}$ \\
 &          & BB & $2.667^{+0.557}_{-0.569}$ & $-1.946^{+0.497}_{-0.595}$ & $-3.232^{+0.338}_{-0.383}$ & $1.417^{+0.007}_{-0.007}$ & $-2.419^{+0.009}_{-0.009}$ & $1.543^{+0.008}_{-0.008}$ & $0.088^{+0.021}_{-0.020}$ \\
 & QJ+WP+Pl & EE & $7.889^{+0.332}_{-0.335}$ & $-2.816^{+0.094}_{-0.095}$ & $-3.130^{+0.042}_{-0.042}$ & $1.889^{+0.011}_{-0.011}$ & $-2.111^{+0.009}_{-0.009}$ & $1.533^{+0.008}_{-0.008}$ & $0.051^{+0.009}_{-0.009}$ \\
 &          & BB & $1.266^{+0.371}_{-0.355}$ & $-3.861^{+0.570}_{-0.697}$ & $-2.986^{+0.136}_{-0.135}$ & $1.416^{+0.007}_{-0.007}$ & $-2.418^{+0.009}_{-0.009}$ & $1.542^{+0.008}_{-0.008}$ & $0.062^{+0.018}_{-0.017}$ \\
\midrule
\multirow{4}{*}{$|b|>20^\circ$}
 & WP+Pl    & EE & $7.902^{+0.720}_{-0.716}$ & $-2.957^{+0.203}_{-0.216}$ & $-3.260^{+0.145}_{-0.151}$ & $1.362^{+0.010}_{-0.010}$ & $-2.205^{+0.011}_{-0.011}$ & $1.487^{+0.011}_{-0.011}$ & $0.056^{+0.013}_{-0.013}$ \\
 &          & BB & $2.675^{+0.588}_{-0.613}$ & $-1.663^{+0.543}_{-0.661}$ & $-3.055^{+0.350}_{-0.403}$ & $0.949^{+0.006}_{-0.006}$ & $-2.335^{+0.011}_{-0.011}$ & $1.540^{+0.011}_{-0.010}$ & $0.079^{+0.025}_{-0.023}$ \\
 & QJ+WP+Pl & EE & $7.218^{+0.349}_{-0.347}$ & $-3.010^{+0.108}_{-0.111}$ & $-3.134^{+0.044}_{-0.044}$ & $1.362^{+0.010}_{-0.010}$ & $-2.205^{+0.011}_{-0.011}$ & $1.486^{+0.011}_{-0.011}$ & $0.033^{+0.010}_{-0.010}$ \\
 &          & BB & $1.369^{+0.408}_{-0.401}$ & $-3.561^{+0.601}_{-0.744}$ & $-2.918^{+0.152}_{-0.152}$ & $0.949^{+0.006}_{-0.006}$ & $-2.335^{+0.011}_{-0.011}$ & $1.539^{+0.010}_{-0.010}$ & $0.066^{+0.021}_{-0.020}$ \\
\bottomrule
\end{tabular}}
\end{adjustbox}
\end{table*}

% ----------------------------------------------------
% Individual tables
% ----------------------------------------------------

% \begin{table*}[h]
\centering
\footnotesize
\setlength{\tabcolsep}{4pt}
\caption{Posterior constraints (median and $68\%$ credible intervals) on the synchrotron and dust amplitudes, spectral indices, and synchrotron-dust correlation coefficient $\rho$, from a global power-law fit to the EE and BB polarization auto- and cross-spectra of the 14-band QUIJOTE+WMAP+Planck dataset over the multipole range $30 \leq \ell \leq 199$, using the $|b|>10^\circ$ Galactic mask.}
\label{tab:fit_galcut10_all_pairs_ell30-200}
\scalebox{0.95}{
\begin{tabular}{llcccccccc}
\toprule
Data & Mode & $A_s\,[10^{-3}\,\mu\mathrm{K}^2]$ & $\alpha_s$ & $\beta_s$ & $A_d\,[10^{-3}\,\mu\mathrm{K}^2]$ & $\alpha_d$ & $\beta_d$ & $\rho$ & $\chi^2_\mathrm{red}$ \\
\midrule
WP+Pl & EE & $9.451^{+0.649}_{-0.644}$ & $-2.845^{+0.146}_{-0.156}$ & $-3.189^{+0.116}_{-0.120}$ & $2.999^{+0.012}_{-0.012}$ & $-2.244^{+0.006}_{-0.006}$ & $1.523^{+0.006}_{-0.006}$ & $0.090^{+0.011}_{-0.011}$ & $3.99$ \\
WP+Pl & BB & $2.589^{+0.534}_{-0.530}$ & $-2.380^{+0.443}_{-0.506}$ & $-3.402^{+0.295}_{-0.325}$ & $2.336^{+0.008}_{-0.008}$ & $-2.163^{+0.006}_{-0.006}$ & $1.536^{+0.005}_{-0.005}$ & $0.125^{+0.020}_{-0.019}$ & $2.93$ \\
\midrule
QJ+WP+Pl & EE & $8.806^{+0.325}_{-0.328}$ & $-2.949^{+0.079}_{-0.083}$ & $-3.148^{+0.036}_{-0.036}$ & $2.999^{+0.012}_{-0.012}$ & $-2.244^{+0.006}_{-0.006}$ & $1.523^{+0.006}_{-0.006}$ & $0.088^{+0.008}_{-0.008}$ & $3.27$ \\
QJ+WP+Pl & BB & $1.266^{+0.296}_{-0.284}$ & $-4.074^{+0.449}_{-0.535}$ & $-3.065^{+0.109}_{-0.110}$ & $2.336^{+0.008}_{-0.008}$ & $-2.163^{+0.006}_{-0.006}$ & $1.535^{+0.005}_{-0.005}$ & $0.103^{+0.016}_{-0.015}$ & $2.51$ \\
\bottomrule
\end{tabular}}
\end{table*}


% \begin{table*}[h]
\centering
\footnotesize
\setlength{\tabcolsep}{4pt}
\caption{Posterior constraints (median and $68\%$ credible intervals) on the synchrotron and dust amplitudes, spectral indices, and synchrotron-dust correlation coefficient $\rho$, from a global power-law fit to the EE and BB polarization auto- and cross-spectra of the 14-band QUIJOTE+WMAP+Planck dataset over the multipole range $30 \leq \ell \leq 199$, using the $|b|>15^\circ$ Galactic mask.}
\label{tab:fit_galcut15_all_pairs_ell30-200}
\scalebox{0.95}{
\begin{tabular}{llcccccccc}
\toprule
Data & Mode & $A_s\,[10^{-3}\,\mu\mathrm{K}^2]$ & $\alpha_s$ & $\beta_s$ & $A_d\,[10^{-3}\,\mu\mathrm{K}^2]$ & $\alpha_d$ & $\beta_d$ & $\rho$ & $\chi^2_\mathrm{red}$ \\
\midrule
WP+Pl & EE & $8.697^{+0.678}_{-0.677}$ & $-2.695^{+0.169}_{-0.181}$ & $-3.234^{+0.138}_{-0.143}$ & $1.889^{+0.011}_{-0.011}$ & $-2.112^{+0.009}_{-0.009}$ & $1.533^{+0.008}_{-0.008}$ & $0.066^{+0.012}_{-0.012}$ & $2.22$ \\
WP+Pl & BB & $2.667^{+0.557}_{-0.569}$ & $-1.946^{+0.497}_{-0.595}$ & $-3.232^{+0.338}_{-0.383}$ & $1.417^{+0.007}_{-0.007}$ & $-2.419^{+0.009}_{-0.009}$ & $1.543^{+0.008}_{-0.008}$ & $0.088^{+0.021}_{-0.020}$ & $1.63$ \\
\midrule
QJ+WP+Pl & EE & $7.889^{+0.332}_{-0.335}$ & $-2.816^{+0.094}_{-0.095}$ & $-3.130^{+0.042}_{-0.042}$ & $1.889^{+0.011}_{-0.011}$ & $-2.111^{+0.009}_{-0.009}$ & $1.533^{+0.008}_{-0.008}$ & $0.051^{+0.009}_{-0.009}$ & $1.95$ \\
QJ+WP+Pl & BB & $1.266^{+0.371}_{-0.355}$ & $-3.861^{+0.570}_{-0.697}$ & $-2.986^{+0.136}_{-0.135}$ & $1.416^{+0.007}_{-0.007}$ & $-2.418^{+0.009}_{-0.009}$ & $1.542^{+0.008}_{-0.008}$ & $0.062^{+0.018}_{-0.017}$ & $1.57$ \\
\bottomrule
\end{tabular}}
\end{table*}


% \begin{table*}[h]
\centering
\footnotesize
\setlength{\tabcolsep}{4pt}
\caption{Posterior constraints (median and $68\%$ credible intervals) on the synchrotron and dust amplitudes, spectral indices, and synchrotron-dust correlation coefficient $\rho$, from a global power-law fit to the EE and BB polarization auto- and cross-spectra of the 14-band QUIJOTE+WMAP+Planck dataset over the multipole range $30 \leq \ell \leq 199$, using the $|b|>20^\circ$ Galactic mask.}
\label{tab:fit_galcut20_all_pairs_ell30-200}
\scalebox{0.95}{
\begin{tabular}{llcccccccc}
\toprule
Data & Mode & $A_s\,[10^{-3}\,\mu\mathrm{K}^2]$ & $\alpha_s$ & $\beta_s$ & $A_d\,[10^{-3}\,\mu\mathrm{K}^2]$ & $\alpha_d$ & $\beta_d$ & $\rho$ & $\chi^2_\mathrm{red}$ \\
\midrule
WP+Pl & EE & $7.902^{+0.720}_{-0.716}$ & $-2.957^{+0.203}_{-0.216}$ & $-3.260^{+0.145}_{-0.151}$ & $1.362^{+0.010}_{-0.010}$ & $-2.205^{+0.011}_{-0.011}$ & $1.487^{+0.011}_{-0.011}$ & $0.056^{+0.013}_{-0.013}$ & $1.83$ \\
WP+Pl & BB & $2.675^{+0.588}_{-0.613}$ & $-1.663^{+0.543}_{-0.661}$ & $-3.055^{+0.350}_{-0.403}$ & $0.949^{+0.006}_{-0.006}$ & $-2.335^{+0.011}_{-0.011}$ & $1.540^{+0.011}_{-0.010}$ & $0.079^{+0.025}_{-0.023}$ & $1.47$ \\
\midrule
QJ+WP+Pl & EE & $7.218^{+0.349}_{-0.347}$ & $-3.010^{+0.108}_{-0.111}$ & $-3.134^{+0.044}_{-0.044}$ & $1.362^{+0.010}_{-0.010}$ & $-2.205^{+0.011}_{-0.011}$ & $1.486^{+0.011}_{-0.011}$ & $0.033^{+0.010}_{-0.010}$ & $1.69$ \\
QJ+WP+Pl & BB & $1.369^{+0.408}_{-0.401}$ & $-3.561^{+0.601}_{-0.744}$ & $-2.918^{+0.152}_{-0.152}$ & $0.949^{+0.006}_{-0.006}$ & $-2.335^{+0.011}_{-0.011}$ & $1.539^{+0.010}_{-0.010}$ & $0.066^{+0.021}_{-0.020}$ & $1.42$ \\
\bottomrule
\end{tabular}}
\end{table*}


% ----------------------------------------------------

\subsection{Fitting only cross-spectra}

\textcolor{red}{RESULTS CROSS-SPECTRA ONLY -- $30 \leq \ell \leq 199$}

\begin{table*}[h]
\centering
\footnotesize
\setlength{\tabcolsep}{4pt}
\caption{Posterior constraints (median and $68\%$ credible intervals) on the synchrotron and dust amplitudes, spectral indices, and synchrotron-dust correlation coefficient $\rho$, from a global power-law fit to the EE and BB polarization cross-spectra of the 14-band QUIJOTE+WMAP+Planck dataset over the multipole range $30 \leq \ell \leq 199$, for the three Galactic masks.}
\label{tab:fit_all_masks_cross_only_ell30-200}
\begin{adjustbox}{center}
\scalebox{1}{
\begin{tabular}{llllcccccccc}
\toprule
Mask & Dataset & Mode & $A_s\,[10^{-3}\,\mu\mathrm{K}^2]$ & $\alpha_s$ & $\beta_s$ & $A_d\,[10^{-3}\,\mu\mathrm{K}^2]$ & $\alpha_d$ & $\beta_d$ & $\rho$ \\
\midrule
\multirow{4}{*}{$|b|>10^\circ$}
 & WP+Pl    & EE & $11.794^{+1.193}_{-1.126}$ & $-2.661^{+0.164}_{-0.172}$ & $-3.550^{+0.203}_{-0.225}$ & $2.962^{+0.022}_{-0.022}$ & $-2.240^{+0.008}_{-0.008}$ & $1.512^{+0.010}_{-0.009}$ & $0.088^{+0.010}_{-0.010}$ \\
 &          & BB & $3.854^{+1.018}_{-0.885}$ & $-2.075^{+0.402}_{-0.447}$ & $-3.981^{+0.466}_{-0.553}$ & $2.372^{+0.016}_{-0.016}$ & $-2.160^{+0.008}_{-0.008}$ & $1.561^{+0.009}_{-0.009}$ & $0.113^{+0.019}_{-0.017}$ \\
 & QJ+WP+Pl & EE & $9.027^{+0.360}_{-0.362}$ & $-2.868^{+0.089}_{-0.090}$ & $-3.170^{+0.046}_{-0.047}$ & $2.967^{+0.022}_{-0.022}$ & $-2.240^{+0.008}_{-0.008}$ & $1.515^{+0.010}_{-0.009}$ & $0.087^{+0.008}_{-0.008}$ \\
 &          & BB & $1.642^{+0.319}_{-0.321}$ & $-3.482^{+0.406}_{-0.470}$ & $-2.891^{+0.120}_{-0.119}$ & $2.373^{+0.016}_{-0.016}$ & $-2.159^{+0.008}_{-0.008}$ & $1.564^{+0.009}_{-0.009}$ & $0.106^{+0.016}_{-0.015}$ \\
\midrule
\multirow{4}{*}{$|b|>15^\circ$}
 & WP+Pl    & EE & $11.700^{+1.432}_{-1.278}$ & $-2.396^{+0.177}_{-0.190}$ & $-3.737^{+0.265}_{-0.306}$ & $1.872^{+0.019}_{-0.019}$ & $-2.110^{+0.011}_{-0.011}$ & $1.525^{+0.014}_{-0.013}$ & $0.066^{+0.011}_{-0.011}$ \\
 &          & BB & $4.234^{+1.211}_{-1.041}$ & $-1.638^{+0.464}_{-0.515}$ & $-4.233^{+0.637}_{-0.770}$ & $1.422^{+0.014}_{-0.013}$ & $-2.420^{+0.012}_{-0.012}$ & $1.551^{+0.013}_{-0.013}$ & $0.081^{+0.019}_{-0.017}$ \\
 & QJ+WP+Pl & EE & $8.104^{+0.364}_{-0.363}$ & $-2.706^{+0.100}_{-0.104}$ & $-3.151^{+0.055}_{-0.055}$ & $1.873^{+0.019}_{-0.019}$ & $-2.110^{+0.011}_{-0.011}$ & $1.526^{+0.013}_{-0.013}$ & $0.051^{+0.009}_{-0.009}$ \\
 &          & BB & $1.685^{+0.367}_{-0.387}$ & $-3.171^{+0.480}_{-0.594}$ & $-2.857^{+0.138}_{-0.145}$ & $1.423^{+0.014}_{-0.014}$ & $-2.420^{+0.012}_{-0.012}$ & $1.553^{+0.013}_{-0.013}$ & $0.063^{+0.017}_{-0.017}$ \\
\midrule
\multirow{4}{*}{$|b|>20^\circ$}
 & WP+Pl    & EE & $11.296^{+1.558}_{-1.375}$ & $-2.647^{+0.211}_{-0.228}$ & $-3.901^{+0.304}_{-0.352}$ & $1.334^{+0.017}_{-0.017}$ & $-2.211^{+0.015}_{-0.014}$ & $1.460^{+0.017}_{-0.017}$ & $0.055^{+0.012}_{-0.012}$ \\
 &          & BB & $4.148^{+1.380}_{-1.156}$ & $-1.387^{+0.532}_{-0.621}$ & $-4.149^{+0.743}_{-0.914}$ & $0.950^{+0.012}_{-0.012}$ & $-2.337^{+0.015}_{-0.015}$ & $1.543^{+0.018}_{-0.018}$ & $0.075^{+0.021}_{-0.019}$ \\
 & QJ+WP+Pl & EE & $7.481^{+0.376}_{-0.378}$ & $-2.880^{+0.115}_{-0.120}$ & $-3.154^{+0.057}_{-0.057}$ & $1.334^{+0.017}_{-0.017}$ & $-2.210^{+0.014}_{-0.014}$ & $1.461^{+0.017}_{-0.017}$ & $0.032^{+0.010}_{-0.010}$ \\
 &          & BB & $1.760^{+0.384}_{-0.407}$ & $-2.906^{+0.506}_{-0.613}$ & $-2.827^{+0.154}_{-0.161}$ & $0.952^{+0.012}_{-0.012}$ & $-2.337^{+0.016}_{-0.016}$ & $1.547^{+0.018}_{-0.018}$ & $0.065^{+0.020}_{-0.020}$ \\
\bottomrule
\end{tabular}}
\end{adjustbox}
\end{table*}

% ----------------------------------------------------
% Individual tables
% ----------------------------------------------------

% \begin{table*}[h]
\centering
\scriptsize
\setlength{\tabcolsep}{4pt}
\caption{Posterior constraints (median and $68\%$ credible intervals) on the synchrotron and dust amplitudes and spectral indices, and on the synchrotron--dust correlation coefficient $\rho$, from fits to the $EE$ and $BB$ spectra using the 10º mask (multipole range $\ell=30-200$).}
\label{tab:fit_cross_only_ell30-200}
\begin{tabular}{llcccccccc}
\toprule
Data & Mode & $A_s\,[10^{-3}\,\mu\mathrm{K}^2]$ & $\alpha_s$ & $\beta_s$ & $A_d\,[10^{-3}\,\mu\mathrm{K}^2]$ & $\alpha_d$ & $\beta_d$ & $\rho$ & $\chi^2_\mathrm{red}$ \\
\midrule
WMAP+Planck & EE & $11.582^{+1.179}_{-1.113}$ & $-2.664^{+0.162}_{-0.172}$ & $-3.492^{+0.200}_{-0.222}$ & $3.795^{+0.032}_{-0.031}$ & $-2.239^{+0.008}_{-0.008}$ & $1.587^{+0.010}_{-0.010}$ & $0.081^{+0.009}_{-0.009}$ & $3.15$ \\
WMAP+Planck & BB & $3.615^{+0.991}_{-0.853}$ & $-2.125^{+0.399}_{-0.456}$ & $-3.814^{+0.459}_{-0.539}$ & $3.050^{+0.024}_{-0.024}$ & $-2.160^{+0.008}_{-0.008}$ & $1.638^{+0.010}_{-0.010}$ & $0.107^{+0.018}_{-0.016}$ & $2.28$ \\
QUIJOTE+WMAP+Planck & EE & $9.078^{+0.362}_{-0.364}$ & $-2.866^{+0.088}_{-0.092}$ & $-3.162^{+0.047}_{-0.046}$ & $3.801^{+0.032}_{-0.032}$ & $-2.239^{+0.008}_{-0.008}$ & $1.589^{+0.010}_{-0.010}$ & $0.080^{+0.007}_{-0.007}$ & $2.62$ \\
QUIJOTE+WMAP+Planck & BB & $1.684^{+0.313}_{-0.312}$ & $-3.450^{+0.390}_{-0.444}$ & $-2.859^{+0.118}_{-0.120}$ & $3.052^{+0.024}_{-0.023}$ & $-2.159^{+0.008}_{-0.008}$ & $1.640^{+0.010}_{-0.010}$ & $0.099^{+0.015}_{-0.014}$ & $2.01$ \\
\bottomrule
\end{tabular}
\end{table*}

% \begin{table*}[h]
\centering
\footnotesize
\setlength{\tabcolsep}{4pt}
\caption{Posterior constraints (median and $68\%$ credible intervals) on the synchrotron and dust amplitudes, spectral indices, and synchrotron-dust correlation coefficient $\rho$, from a global power-law fit to the EE and BB polarization cross-spectra of the 14-band QUIJOTE+WMAP+Planck dataset over the multipole range $30 \leq \ell \leq 199$, using the $|b|>15^\circ$ Galactic mask.}
\label{tab:fit_galcut15_cross_only_ell30-200}
\scalebox{0.95}{
\begin{tabular}{llcccccccc}
\toprule
Data & Mode & $A_s\,[10^{-3}\,\mu\mathrm{K}^2]$ & $\alpha_s$ & $\beta_s$ & $A_d\,[10^{-3}\,\mu\mathrm{K}^2]$ & $\alpha_d$ & $\beta_d$ & $\rho$ & $\chi^2_\mathrm{red}$ \\
\midrule
WP+Pl & EE & $11.700^{+1.432}_{-1.278}$ & $-2.396^{+0.177}_{-0.190}$ & $-3.737^{+0.265}_{-0.306}$ & $1.872^{+0.019}_{-0.019}$ & $-2.110^{+0.011}_{-0.011}$ & $1.525^{+0.014}_{-0.013}$ & $0.066^{+0.011}_{-0.011}$ & $1.92$ \\
WP+Pl & BB & $4.234^{+1.211}_{-1.041}$ & $-1.638^{+0.464}_{-0.515}$ & $-4.233^{+0.637}_{-0.770}$ & $1.422^{+0.014}_{-0.013}$ & $-2.420^{+0.012}_{-0.012}$ & $1.551^{+0.013}_{-0.013}$ & $0.081^{+0.019}_{-0.017}$ & $1.48$ \\
\midrule
QJ+WP+Pl & EE & $8.104^{+0.364}_{-0.363}$ & $-2.706^{+0.100}_{-0.104}$ & $-3.151^{+0.055}_{-0.055}$ & $1.873^{+0.019}_{-0.019}$ & $-2.110^{+0.011}_{-0.011}$ & $1.526^{+0.013}_{-0.013}$ & $0.051^{+0.009}_{-0.009}$ & $1.73$ \\
QJ+WP+Pl & BB & $1.685^{+0.367}_{-0.387}$ & $-3.171^{+0.480}_{-0.594}$ & $-2.857^{+0.138}_{-0.145}$ & $1.423^{+0.014}_{-0.014}$ & $-2.420^{+0.012}_{-0.012}$ & $1.553^{+0.013}_{-0.013}$ & $0.063^{+0.017}_{-0.017}$ & $1.45$ \\
\bottomrule
\end{tabular}}
\end{table*}


% \begin{table*}[h]
\centering
\footnotesize
\setlength{\tabcolsep}{4pt}
\caption{Posterior constraints (median and $68\%$ credible intervals) on the synchrotron and dust amplitudes, spectral indices, and synchrotron-dust correlation coefficient $\rho$, from a global power-law fit to the EE and BB polarization cross-spectra of the 14-band QUIJOTE+WMAP+Planck dataset over the multipole range $30 \leq \ell \leq 199$, using the $|b|>20^\circ$ Galactic mask.}
\label{tab:fit_galcut20_cross_only_ell30-200}
\scalebox{0.95}{
\begin{tabular}{llcccccccc}
\toprule
Data & Mode & $A_s\,[10^{-3}\,\mu\mathrm{K}^2]$ & $\alpha_s$ & $\beta_s$ & $A_d\,[10^{-3}\,\mu\mathrm{K}^2]$ & $\alpha_d$ & $\beta_d$ & $\rho$ & $\chi^2_\mathrm{red}$ \\
\midrule
WP+Pl & EE & $11.296^{+1.558}_{-1.375}$ & $-2.647^{+0.211}_{-0.228}$ & $-3.901^{+0.304}_{-0.352}$ & $1.334^{+0.017}_{-0.017}$ & $-2.211^{+0.015}_{-0.014}$ & $1.460^{+0.017}_{-0.017}$ & $0.055^{+0.012}_{-0.012}$ & $1.57$ \\
WP+Pl & BB & $4.148^{+1.380}_{-1.156}$ & $-1.387^{+0.532}_{-0.621}$ & $-4.149^{+0.743}_{-0.914}$ & $0.950^{+0.012}_{-0.012}$ & $-2.337^{+0.015}_{-0.015}$ & $1.543^{+0.018}_{-0.018}$ & $0.075^{+0.021}_{-0.019}$ & $1.39$ \\
\midrule
QJ+WP+Pl & EE & $7.481^{+0.376}_{-0.378}$ & $-2.880^{+0.115}_{-0.120}$ & $-3.154^{+0.057}_{-0.057}$ & $1.334^{+0.017}_{-0.017}$ & $-2.210^{+0.014}_{-0.014}$ & $1.461^{+0.017}_{-0.017}$ & $0.032^{+0.010}_{-0.010}$ & $1.52$ \\
QJ+WP+Pl & BB & $1.760^{+0.384}_{-0.407}$ & $-2.906^{+0.506}_{-0.613}$ & $-2.827^{+0.154}_{-0.161}$ & $0.952^{+0.012}_{-0.012}$ & $-2.337^{+0.016}_{-0.016}$ & $1.547^{+0.018}_{-0.018}$ & $0.065^{+0.020}_{-0.020}$ & $1.36$ \\
\bottomrule
\end{tabular}}
\end{table*}


% ----------------------------------------------------
